\subsection{从操作依赖到真实可执行任务的两级展开}
外层变量 $p=(\pi,m,\sigma)$ 只有子图标识、核归属及同核顺序，尚不是可直接执行的周期时间线。它必须经过两个确定性映射：首先是场景规则 $\mathfrak G_s$，依据场景 $s\in\{A,B,C\}$ 把跨子图张量改写成必要的 COPY、同步和 Task 释放；随后由核内启发式 $\mathfrak K$ 为各 Task 建立 Pipe 顺序与容量安全的换入换出；最后事件仿真器 $\mathfrak E_s$ 给出性能。可把评估写成复合函数
\begin{equation}\label{eq:composition-model}
(T_s(p),D_s(p),\mathcal L_s(p))
=\mathfrak E_s\!\left(\mathfrak K\bigl(\mathfrak G_s(G,p)\bigr)\right),
\end{equation}
其中 $\mathcal L_s(p)$ 为时间线与资源账本。提出一个新划分时，$\mathfrak G_s$ 的输出可能有不同数量的 COPY，$\mathfrak K$ 的换出决定亦可能变化；因此 $T_s(p)$ 不是一个关于子图边界的线性函数。场景 C 在相同图展开后又增加状态相关的 FIFO 路由，故事件映射也随历史轨迹变化。这一复合结构是模型与算法分工的根本：解析式负责刻画合法性、建立下界和排序候选；完整事件评估负责判断候选的真实目标值。

对固定 Task $r$，定义其操作事件集合 $E_r$，每个事件 $e$ 有一个所属 Pipe、前驱集合及可能的张量写入或释放动作。如果把全部事件的持续时间直接相加，得到的是\emph{串行上界式直觉}，并非真实时长。例如 Cube 上的事件 $a$ 与 Vector 上的事件 $b$ 若无相互依赖，时间区间可以相交，Task 长度至少为 $\max\{d_a,d_b\}$，不必达到 $d_a+d_b$；若 $b$ 依赖一个正在 DDR 上读取的张量，它只能等待对应读入结束，而另一个无关的 Cube 仍可继续。因而当 DDR 拥塞增加一段 $\delta$ 周期时，Task 的完成时间增量满足的只有 $0\le \Delta T\le\delta$ 这一类\emph{在固定事件图和单处扰动下}的条件性关系；若延迟改变事件顺序，连这些条件也须重新检查。不能把每条搬运的增量一概加到整个 Task 长度。

\subsection{共享带宽的状态方程与整数周期边界}
为说明仿真为何必须事件驱动，设 DDR 搬运 $e$ 的剩余独占工作量为 $w_e(t)\ge0$，发射事件发生于 $a_e$，初值 $w_e(a_e)=\lceil b_e/B_D\rceil$。在无到达和无完成的区间 $[t,t+\Delta)$，若活跃集合大小为 $n_D$，则
\begin{equation}\label{eq:common-ddr-update}
w_e(t+\Delta)=\max\{0,w_e(t)-\Delta/n_D\},
\qquad e\in A_D(t).
\end{equation}
所能安全推进的下一个完成间隔为 $\Delta_D=n_D\min_{e\in A_D(t)}w_e(t)$；外部新任务、核内 Pipe 事件可能更早发射，因此实际推进量应取 $\min\{\Delta_D,\Delta_{\rm next\ arrival},\Delta_{\rm next\ compute}\}$。事件发生后重新计算活跃集合及其服务份额。官方整数周期语义包含取整和同刻事件处理顺序，式\eqref{eq:common-ddr-update}是解释该机制的连续服务形式；没有复刻同刻顺序的自写代理，不能声称和官方完全等价。该表述也说明“每条边按 $b_e/60$ 加一项”有两类误差：其一，独占时间忽略已有并发搬运在新流加入后变慢；其二，持续时间不能直接转换成最终输出的增量，因为一部分拥塞藏在其它 Pipe 的计算后面。

若对方案 $p$ 的 DDR 事件集合记为 $E_D(p)$，其总独占工作量定义为
\begin{equation}\label{eq:ddr-exclusive-total}
W_D(p)=\sum_{e\in E_D(p)}\left\lceil\frac{b_e}{B_D}\right\rceil.
\end{equation}
连续公平服务在所有在途传输之间重新分配这份工作。若每条传输的整数周期取整后必须服务至少 $\lceil b_e/B_D\rceil$ 个“单流等效周期”，便有 $T(p)\ge W_D(p)$；但是当官方把取整置于不同层级时，以字节总量给出的 $T(p)\ge B_D(p)/B_D$ 才是更稳妥的容量界。本文在理论式中采用后者，并把独占工作量只当作局部代理输入。这样区分可避免把对官方舍入细节的假设写成不加条件的严格下界。

\subsection{关键路径归因与重叠时间的去重}
给定已经评估的时间线，取最后完成的输出事件 $z$，沿使其开始时刻达到最大值的约束边逆行追溯。若 $z$ 的前一事件完成于 $t_1$、其跨核数据直到 $t_2$ 才释放、其 Pipe 直到 $t_3$ 才空闲，则其发射为 $\max\{t_1,t_2,t_3\}$，真正使其等待的只有达到最大值的约束；另外两个可能同时存在，却不是该次发射的边际原因。多个约束同时达到最大值时存在并列关键边，不能强制把同一等待周期完全分给某一个原因。为避免重叠时间多计，可为每段 $[u,v)$ 定义同时满足的状态指标：DDR 忙碌、相关任务未释放、目标 Pipe 空闲。三类指标会重叠，求和可能大于 $v-u$。只有在逐事件逆向追溯后，把\emph{关键链上唯一或并列的激活约束}计入候选诊断，才能解释改善机会。

这一方法也限定了“减少跨核边”和“减少额外搬运”两类结构指标的用途。若一条边处于非关键分支，即使删去其 $1000$ 周期释放规则，也可能不改变最终完成时间；若某个共享输入读入的流量完全与长时计算重叠，减少其 DDR 字节也可能使 $D$ 降低而 $T$ 不动。反之，减少关键输出前的一个很小搬运，可能使后续多个 Task 提前释放，最终节省远高于该条搬运的独占时间。局部候选选择因此考虑边是否与关键输出的祖先集合相交、是否与 DDR 拥塞窗口相交、改变后是否可能触发另一条路径成为关键路径，而不使用单一的全局边权阈值决定接受。

设当前已实现的关键路径长度为 $T_1$，第二长路径为 $T_2$，且候选只缩短第一条路径 $\delta$，不改变其它路径和资源竞争。则新完成时间为 $T'=\max\{T_1-\delta,T_2\}$，边际改善为
\begin{equation}\label{eq:critical-path-saturation}
T_1-T'=\min\{\delta,T_1-T_2\}.
\end{equation}
这是固定路径集合下的条件性推导，说明单个瓶颈的收益会在第二路径接管时饱和。真实方案改变会重构路径，式\eqref{eq:critical-path-saturation}仅用作排序和解释，不作为官方评价结果。它也解释为何反复围绕一条旧关键链做开环搜索，后期容易出现大量完全持平的候选；继续增大区域未必有效，可能只增加评价机会成本。

\subsection{候选排序的可证范围与公平预算}
一个候选池 $\mathcal C=\{p_1,\ldots,p_M\}$ 中，代理给出 $\widehat T_j$，官方评价给出 $T_j$。若代理只能评价所有 $M$ 个、官方最多评价 $B<M$ 个，选择器的实际损失取决于前 $B$ 个覆盖了多少真优候选，而不是代理均方误差本身。令 $j^*=\arg\min_j T_j$，代理按升序排列索引为 $\widehat j_1,\ldots,\widehat j_M$，可定义候选池内选择遗憾
\begin{equation}\label{eq:selection-regret}
R_B=\min_{1\le b\le B}T_{\widehat j_b}-T_{j^*}\ge0.
\end{equation}
当 $R_B=0$ 时，前 $B$ 覆盖了该池的官方最优；但这\emph{不表示}全局最优已经找到。若同一指纹的方案重复占用多个名额，名义预算 $B$ 与真正不同的 $B_{\rm unique}$ 不同，必须按唯一完整方案数登记。冻结同一核内准备与成本状态后，才能把多个候选的代理名次和官方名次相配对，分析误差来自 Task 释放、Pipe 路径还是共享 DDR；若在候选间随意重算准备状态，相关系数和遗憾的解释基础也消失。

公平对照使用相同的图、核数、代码版本、随机种子、官方评价次数和候选保底。若新增一个机制挤掉原来有效的候选，它的表观收益包含预算替代效应；若只给新机制额外评价机会，则结果又混入预算扩张效应。定义 $\mathcal B$ 为固定官方评价机会，基础搜索分配 $B_0$，结构候选分配 $B_1$，局部修复分配 $B_2$，需满足 $B_0+B_1+B_2\le\mathcal B$。在某一邻域连续无效时可以把尚未使用的 $B_1$ 或 $B_2$ 转给另一个，但必须保留基础搜索下限，且在日志记录实际采用的次数；否则同样的名义预算可能对应不同的有效探索。本文有限预算优化的目标是高质量地使用这 $\mathcal B$ 次官方调用，而不是把公式下界当成可直接提交的方案。

\subsection{可行性、最优性与求解效率的三重区分}
可行性由唯一覆盖、商图无环、场景同步规则和物理存储约束决定；性能由 $T_s(p)$ 和 $D_s(p)$ 决定；计算效率则由生成并评价 $p$ 的主机耗时决定。这三者在论文中分别检验。一个有效地压缩了 CPU 搜索时间的候选缓存，不一定改变 NPU 执行周期；一个周期很小的代理预测方案，若未经过正式评价与内存检查，不能称为高质量可行结果；一个正式评价优于对照的方案，也不能据此推断全局最优。

对每个固定候选池 $\mathcal C$，官方评价中的最小值 $\min_{p\in\mathcal C}T_s(p)$ 只提供\emph{该池上的}上界，即全局最优 $T_s^*$ 满足 $T_s^*\le\min_{p\in\mathcal C}T_s(p)$。通用计算下界提供相反方向的 $L_0(k)\le T_s^*$。于是合法的、但一般较宽的夹逼关系为
\begin{equation}\label{eq:global-bound-sandwich}
L_0(k)\le T_s^*\le\min_{p\in\mathcal C}T_s(p).
\end{equation}
这个区间明确表明已找到方案的质量边界，也避免把某个方案自身的搬运下界误用为全局下界。当 $L_0(k)$ 很低时，不能从宽区间推断还剩多少可实现提升；当某一大图距 $L_0(k)$ 很远时，也不能断言差距都可消除，因为同步、内存与带宽尚未进入 $L_0$。因此本文以同口径官方结果、负例和结构诊断来评价启发式，而不虚构最优性证明。
